\documentclass{article}

\usepackage[spanish]{babel}
\usepackage{amsmath,amssymb}

\title{Notación Asintótica}
\author{Benjamín Rivas Beltrán}
\date{\today}

\begin{document}

\maketitle

\section{Notación Theta}

Sea $g(n)$ una función que describe el tiempo de ejecución de un algoritmo en función del tamaño de la entrada $n$. Definimos la notación Theta ($\Theta$) de la siguiente manera:
\[
\Theta(g(n)) = \{ f(n) : \exists c_1, c_2 > 0, n_0 \geq 0 \text{ tal que } c_1 g(n) \leq f(n) \leq c_2 g(n) \text{ para todo } n \geq n_0 \}
\]

\subsection{Propiedades de la notación Theta}
\begin{itemize}
  \item $f(n) \in \Theta(g(n))$ implica que $f(n)$ y $g(n)$ crecen al mismo ritmo asintóticamente.
  \item La notación Theta está pensado para funciones monótonas y no negativas.
  \item Existe la equivalencia:
  \[
  f(n) \in \Theta(g(n)) \iff \lim_{n \to \infty} \frac{f(n)}{g(n)} = c, \text{ donde } 0 < c < \infty
  \]
\end{itemize}
\end{document}